% Faithful conversion of ages-i-sprint-3; original XML and source characters retained.
\hypertarget{__RefHeading___Toc4749_3118176633}{}
\refstepcounter{subsection}
\subsection*{2.3.4 Sprint 3}
\label{unit:ages-i-sprint-3}
\addcontentsline{toc}{subsection}{2.3.4 Sprint 3}

% text-p-0557
\noindent \originaltab{}A Sprint 3 foi a mais rápida entre todas as Sprints que já tivemos. Devido a isso, reunimos todos os esforços possíveis para atacar as User Stories da Sprint atual e das anteriores que não havíamos entregue por completo. Mesmo com alguns desfalques em nossa equipe, conseguimos entregar apresentar funcionalidades de grande importância para os clientes. Cabe aqui apontar que três das seis User Stories atacadas foram entregues principalmente por mim.

% text-p-0558
\noindent \originaltab{}Iniciei a Sprint 3 realizando diversas modificações em minha task dos endpoints que adicionam e removem os suplementos do estoque do lojista. Essa tarefa tinha um prazo de entrega para a Sprint 2. Porém, por conta das enchentes, meus colegas de equipe não puderam fazer a review do código e mesclar com o projeto. Devido a isso, decidi, juntamente do Ages IV Arthur Ibarra, revisar a tarefa para entregar nessa Sprint.\originalspace{}

% text-p-0559
\noindent \originaltab{}Haviam alguns erros de compreensão da tarefa por minha parte. Eu havia desenvolvido endpoints capazes de criar e remover suplementos do Banco de Dados. Tais endpoints devem ser disponíveis apenas aos administradores do Easy Choose. Os lojistas só possuem permissão para adicionar e remover de seus estoques os suplementos que já estão disponíveis no Banco de Dados.\originalspace{}

% text-p-0560
\noindent \originaltab{}Identificados os problemas, me juntei aos colegas (AGES II) Matheus Chong e (AGES IV) Arthur Ibarra para realizar um trio programming (conceito que inventei agora). Sob a supervisão de meus companheiros mais experientes, modifiquei as funções dos endpoints e testei as requisições ao Banco de Dados através do software ThunderBolt. Durante o processo, meus colegas me ensinaram mais a fundo sobre requisições HTTP, Repository, Service e Controller, conceitos essenciais de arquitetura de software e desenvolvimento de aplicações web.\originalspace{}

% text-p-0561
\noindent \originaltab{}Após a entrega e merge da tarefa, encontrei outra task de endopoints no Backlog da Sprint. Como eu havia trabalhado com os padrões de design MVC (Model-View-Controller) recentemente, decidi embarcar em mais um desafio. Desta vez, escolhi desenvolver toda a tarefa sem nenhum tipo de auxílio para testar meus conhecimentos. Esta, consistia em criar os endpoints responsáveis por retornar todos os suplementos disponíveis e não disponíveis pelo lojista. Tive alguns percalços, porém concluí a tarefa.

% text-p-0562
\noindent \originaltab{}Para finalizar, meu colega Arthur Ibarra me designou uma última tarefa antes da apresentação da Sprint aos stakehoders, que, segundo ele, era a task perfeita para mim. A tarefa consistia em criar as funções da API e hooks que conectam o Banco de Dados à tela de Lojista. Pude então, realizar o desenvolvimento completo do processo de uma requisição HTTP — desde a criação das funções da API até os hooks que conectam o Front-end ao Back-end.

\begin{reflection}
% text-p-0565
\noindent \originaltab{}Por mais curta que tenha sido, a Sprint 3 me trouxe diversos novos conhecimentos. Nela, manuseei códigos dos outros 2 repositórios que estamos trabalhando. Pude aprender mais a fundo conceitos de MVC, testes de requisições ao Banco de Dados, funções API e hooks. Não tive medo de fazer questionamentos aos meus colegas e fico grato pela disponibilidade e vontade de me ensinar e instruir a como realizar algumas tarefas. Pude conhecer novas tecnologias, softwares e aprimorar minhas habilidades de depuração e validação de código. A experiência de participar da AGES tem sido muito engrandecedora graças à troca de conhecimentos entre o time.\originalspace{}

% text-p-0566
\noindent \originaltab{}Gostaria também de mencionar uma pequena contribuição minha ao meu colega Gustavo Pias, em que o auxiliei a configurar o Docker e o Banco de Dados do projeto. Só pude realizar tal ato graças aos conhecimentos adquiridos nas aulas de Banco de Dados II, ministradas pelo professor Eduardo Arruda. São nos pequenos detalhes que percebemos a enorme importância das cadeiras cursadas para a formação de um profissional de software completo.\originalspace{}

\end{reflection}
